\subsection{Whole Gauss rays in actual principal-open lines}
\label{subsec:line-gauss-paths}

\begin{theorem}[Actual centered Gauss points and radius-zero evaluation]
\label{spectrum:centered-gauss-whole-ray}
\lean{EnvelopeProperty.Berkovich.RelativeSpectrum.centeredGauss,
EnvelopeProperty.Berkovich.RelativeSpectrum.centeredGauss_val_apply,
EnvelopeProperty.Berkovich.RelativeSpectrum.centeredGauss_val_X_sub_C,
EnvelopeProperty.Berkovich.RelativeSpectrum.centeredGauss_zero,
EnvelopeProperty.Berkovich.RelativeSpectrum.continuous_centeredGauss_prod,
EnvelopeProperty.Berkovich.RelativeSpectrum.continuous_centeredGauss,
EnvelopeProperty.Berkovich.RelativeSpectrum.centeredGauss_injective}
\leanok
\uses{spectrum:gauss-point,spectrum:gauss-joint-continuity,spectrum:gauss-zero-radius}
Let $R,A$ be arbitrary commutative rings, $A$ an actual $R$-algebra,
and let $y$ be a whole point of its spectrum above a prescribed
real-valued coefficient seminorm on $R$. For $a\in A$ and every
$r\ge0$, translation $T\mapsto T+a$ followed by the genuine Gauss
extension gives an actual whole point of $A[T]$ above the same
prescribed coefficient seminorm. The value of $T-a$ is exactly $r$.
At $r=0$ the entire point is evaluation at $a$ followed by $y$.
The family is jointly continuous in $y,r$, and its radius path is injective.
\end{theorem}
\begin{proof}
Pull back the actual Gauss point through the genuine polynomial
translation algebra isomorphism, restricted along the original
coefficient algebra. Whole-point pullback retains every original
coefficient value. The translation sends the entire polynomial $T-a$
to $T$, whose Gauss value is exactly $r$. At radius zero the Gauss
value is the actual constant coefficient; the true translated constant
coefficient is $p(a)$ for every original polynomial $p$, so all
polynomial values agree with the genuine evaluation pullback.
Actual Gauss joint continuity and whole-ring restriction give joint
continuity. For a fixed original point, continuity is proved on every
whole polynomial evaluation in the independently defined topology.
The exact $T-a$ value distinguishes every two radii, proving injectivity.
No field, finite-generation, boundedness or positive-radius hypothesis
is used in these constructions and whole-point identities.
\end{proof}

\begin{theorem}[Positive radii retain every nonzero polynomial over a field]
\label{spectrum:centered-gauss-generic-support}
\lean{EnvelopeProperty.Berkovich.RelativeSpectrum.centeredGauss_val_ne_zero,
EnvelopeProperty.Berkovich.RelativeSpectrum.centeredGauss_support_eq_bot}
\leanok
\uses{spectrum:centered-gauss-whole-ray,spectrum:relative-support}
If the coefficient algebra is a field, each positive-radius centered
Gauss point has nonzero value on every nonzero original polynomial.
Its actual scheme support ideal is therefore zero.
\end{theorem}
\begin{proof}
The true translation isomorphism retains nonzeroness of the whole
polynomial. Choose one actual nonzero coefficient of its translated
image. Every valuation on a field has zero value only on zero, so this
coefficient has positive value. A positive radius gives positive value
to its original weighted coefficient. The genuine Gauss coefficient
lower bound then makes the full polynomial value positive. Applying
this to every original polynomial identifies the whole support ideal
with zero; support equality is derived from all function values.
\end{proof}

\begin{theorem}[The whole ray lifts to the actual principal-open line]
\label{spectrum:centered-gauss-principal-open-lift}
\lean{EnvelopeProperty.Berkovich.RelativeSpectrum.eval_val_le_centeredGauss,
EnvelopeProperty.Berkovich.RelativeSpectrum.centeredGauss_val_ne_zero_of_eval,
EnvelopeProperty.Berkovich.RelativeSpectrum.centeredGaussAway,
EnvelopeProperty.Berkovich.RelativeSpectrum.centeredGaussAway_val_algebraMap,
EnvelopeProperty.Berkovich.RelativeSpectrum.centeredGaussAway_val_X_sub_C,
EnvelopeProperty.Berkovich.RelativeSpectrum.continuous_centeredGaussAway,
EnvelopeProperty.Berkovich.RelativeSpectrum.centeredGaussAway_injective}
\leanok
\uses{spectrum:centered-gauss-whole-ray,spectrum:relative-principal-open}
For $g\in A[T]$ suppose the actual evaluation point has nonzero $g$
value. Then the entire centered ray, including radius zero, has nonzero
$g$ value and lifts to every actual localization of $A[T]$ away from $g$.
The coefficient algebra maps must form their actual scalar tower.
Every original polynomial value is preserved, the localized $T-a$
coordinate still has value $r$, and the entire lifted ray is continuous
and injective.
\end{theorem}
\begin{proof}
The true constant coefficient of the translated $g$ is $g(a)$.
Its weighted coefficient lower bound therefore bounds every entire
ray value of $g$ below by its nonzero original evaluation value,
for every nonnegative radius. Use the genuine principal-localization
valuation lift, with this derived nonvanishing condition. Its whole
original-element formula retains every original polynomial value;
in particular it retains the exact centered-coordinate radius.
Every actual localized element has a true fraction representation
$p/g^n$. Its value along the lift is the original ray value of $p$
divided by the original ray value of $g$ to the power $n$.
Both numerator and denominator vary continuously and the denominator
never vanishes. Whole coordinatewise continuity proves continuity
of the entire localization lift. The centered-coordinate identity
again distinguishes all radii. These are actual localizations and
whole valuation extensions, not chosen principal-open certificates.
\end{proof}

\begin{theorem}[Actual evaluation points are not isolated in their line opens]
\label{spectrum:line-evaluation-not-isolated}
\lean{EnvelopeProperty.Berkovich.RelativeSpectrum.exists_centeredGaussAway_ne_in_nhds,
EnvelopeProperty.Berkovich.RelativeSpectrum.not_isOpen_singleton_centeredGaussAway_zero}
\leanok
\uses{spectrum:centered-gauss-principal-open-lift}
Every neighborhood of the entire radius-zero evaluation lift contains
a distinct point of the actual principal-open line at positive radius.
Consequently this whole evaluation point is not isolated. This holds
over the arbitrary coefficient rings and prescribed seminorms above.
\end{theorem}
\begin{proof}
Zero belongs to the closure of the positive radii in the genuine
nonnegative-real topology. Continuity pulls back every neighborhood
of the whole radius-zero lift to a neighborhood of zero.
Intersect its eventual membership with the frequently positive-radius
condition to obtain an actual positive radius whose lift is in the
given neighborhood. Injectivity makes that whole point distinct from
the original radius-zero point. An open singleton would be such a
neighborhood and would contradict this derived distinct-point property.
No maximum principle or subharmonic-function hypothesis is asserted.
\end{proof}

\begin{theorem}[Field evaluation points are approached through generic support]
\label{spectrum:line-evaluation-generic-closure}
\lean{EnvelopeProperty.Berkovich.RelativeSpectrum.centeredGaussAway_val_ne_zero,
EnvelopeProperty.Berkovich.RelativeSpectrum.centeredGaussAway_support_eq_bot,
EnvelopeProperty.Berkovich.RelativeSpectrum.centeredGaussAway_zero_mem_closure_generic}
\leanok
\uses{spectrum:centered-gauss-generic-support,spectrum:centered-gauss-principal-open-lift,spectrum:relative-support}
When the coefficient algebra is a field, each positive-radius localized
point has nonzero value on every nonzero entire localized function
and has zero support ideal. The whole actual evaluation lift is in
the closure of this generic-support locus of the principal-open line.
There is no completeness, algebraic-closure, infinitude, characteristic
or nontrivial-valuation assumption.
\end{theorem}
\begin{proof}
Represent an original nonzero localized function by a true fraction
$p/g^n$. Its numerator is nonzero, since a zero numerator would make
the original fraction zero. The preceding positive-radius field
Gauss theorem makes its entire numerator value nonzero, and the
derived original denominator value is nonzero. The fraction value
is thus nonzero. Applying this to every entire localized function
gives the zero support ideal.

Continuity carries the closure of positive radii into the closure
of their actual point image. Every point of this image has the just
proved generic support, so monotonicity of closure puts the whole
radius-zero evaluation lift in the asserted locus closure.
This is density at the constructed whole evaluation point, derived
in the valuation topology. It does not infer native density from
scheme support density, assert density on every algebraic curve,
or prove a subharmonic maximum principle.
\end{proof}
